% =====================================================================
% TEMPLATE: ARTIGO / NOTAS / LISTA DE EXERCÍCIOS
% Documento curto em article, com os mesmos pacotes e ambientes do TCC.
% Compilar com: latexmk -xelatex -outdir=../build main.tex
% =====================================================================
\documentclass[12pt, a4paper]{article}

% ---------- Configuração dos módulos comuns ----------
\newcommand{\comumNumeracao}{section}        % Teorema 1.1, 1.2, ...
\newcommand{\comumBibliografia}{referencias.bib}
\newcommand{\comumContadorUnico}{}           % Def. 1.1, Teo. 1.2, Lema 1.3, ...

\input{comum/packages}

% Artigo não usa cabeçalho com régua: só o número da página no rodapé.
\pagestyle{plain}

% ---------- Dados do documento ----------
\title{Título do Artigo}
\author{%
  Nome do Autor\thanks{Unespar, \textit{campus} de Campo Mourão.}
  \and
  Nome do Coautor%
}
\date{\today}

\hypersetup{pdftitle={Título do Artigo}, pdfauthor={Nome do Autor}}

\begin{document}

\maketitle

\begin{abstract}
  Resumo do trabalho (entre 100 e 250 palavras).

  \medskip
  \noindent\textbf{Palavras-chave:} palavra 1; palavra 2; palavra 3.
\end{abstract}

\section{Introdução}
\label{sec:introducao}

Texto da introdução \parencite{doCarmo1992}.

\section{Desenvolvimento}
\label{sec:desenvolvimento}

\begin{definition}
  \label{def:exemplo}
  Definição de exemplo.
\end{definition}

\begin{theorem}
  \label{teo:exemplo}
  Pela \cref{def:exemplo}, vale \ldots
\end{theorem}

\begin{proof}
  Demonstração.
\end{proof}

\section{Considerações finais}

Conforme o \Cref{teo:exemplo} \ldots

\printbibliography[title={Referências}]

\end{document}
